% Authored. Numbers come from generated/macros.tex.

\section{Conclusion}
\label{sec:conclusion}

Portugal's general-government balance, decomposed across institutional subsectors
over \PanelYears\ years, has a composition that is asymmetric and remarkably stable.
The Central Government balance is negative in every one of the
\CentralNegativeYears\ observations. The Social Security balance is positive in
\SsfPositiveYears, with a longest unbroken run of \SsfLongestPositiveRun\ years.
Regional and Local Government is the only subsector that changes sign with any
regularity, and it is small in both directions. The aggregate tracks Central
Government at a correlation of \AggregateCentralCorrelation. The aggregate is positive
in \NumPositiveYears\ years, and in each of them the combined non-Social-Security
balance is negative --- from which, given the identity, the rest of that year's
composition follows.

Two results cut against readings that the aggregate series invites.

The first concerns Central Government. Its balance is negative throughout the panel,
but its primary balance is positive in \PrimaryPositiveYears\ of the
\PrimaryObservedYears\ years for which detailed accounts exist, from
\PrimaryPositiveFirst\ to \PrimaryPositiveLast. A persistently negative B.9 does not
imply a persistently negative primary balance, and interest --- averaging
\InterestMeanPct\% of GDP and peaking at \InterestPeakPct\% --- is the arithmetic that
separates them.

The second concerns Social Security. The ESA 2010 Social Security Funds sector and
the CFP budget-execution systems are different accounting objects whose balances
differ in every one of the \SsfBoundaryYears\ years they overlap, by between
\SsfBoundaryMin\ and \SsfBoundaryMax\ million euro. Quoting one where the other
belongs introduces an error of that order with no signal that it has occurred.

\subsection{What this paper does not establish}

The decomposition is an accounting identity. It reallocates a recorded quantity and
estimates nothing, and three inferences do not follow from it.

It does not license a counterfactual. That a positive Social Security balance exceeds
a negative non-Social-Security balance in a given year says those two published
numbers stand in a particular ratio. It says nothing about what the aggregate would
have been under a different institutional arrangement, because that counterfactual
would move contribution rates, entitlements and transfers together and the identity
models none of them.

It does not attribute responsibility. A subsector accounting arithmetically for most
of a movement has not been shown to have produced it.

It does not speak to sustainability. Neither the headline nor the primary balance
determines the debt path, which also depends on stock-flow adjustments that are the
larger term in several years of these data.

Two further limits are properties of the data rather than of the method. Magnitudes
are regime-specific and cannot be pooled across the 1995 splice; the pooled aggregate
mean of \GgMeanPooled\% of GDP describes neither the \GgMeanHistorical\% of
\PanelStart--1994 nor the \GgMeanModern\% of 1995--\PanelEnd. And
\ProvisionalYears\ are provisional at source: they are the years discussed in most
detail here and the years most likely to be restated by a later vintage.

\subsection{Extensions}

Three extensions follow naturally, in ascending order of what they would add.

A wage-bill specification would give the Social Security results an economic
mechanism. Equation~\eqref{eq:ssfchange} establishes that contributions dominate the
recent balance changes; relating the change in contributions to the change in the
aggregate wage bill $W_t = N_t \bar{w}_t$ would take that from an accounting location
to an economic one.

A finer decomposition of revenue and expenditure would sharpen the episode attribution
of Table~\ref{tab:hierarchy}, which uses subsector totals and therefore cannot see
composition shifts within them.

Most valuable would be a European benchmark. ESA 2010 requires the same subsector
breakdown from every member state, and Eurostat and the ECB publish it
\citep{eurostat_gov_main,ecb_gfs}. Defining, for each country and year,
$N_{c,t} = \Bc_{c,t} + \Brl_{c,t}$ and, where the sign condition holds,
$O_{c,t} = \Bssf_{c,t}/|N_{c,t}|$ would place Portugal in a distribution and answer
the question this paper cannot: whether the composition documented here is unusual
or ordinary. As a single-country longitudinal study, nothing above can distinguish a
Portuguese peculiarity from a general feature of European fiscal accounting.
